% Source: GLM23v3, REsubmission.tex, lines 1695--1777.
% Scope: exact sector exchange and both actual periodic endpoints.
% Checked against native CI run 37472300569, build job 112298681251.

\subsection{Exchanging the two periodic endpoints}
\label{sec:mpo_joint_periodic_endpoint_swap}

The second endpoint follows from the first by exchanging the two sector
labels. This exchange acts on the shared endpoint alphabets and on every
rectangular alphabet at once, so all overlaps between block columns are
preserved.

\begin{definition}[Physical and virtual sector exchange]
    \label{def:mpo_joint_endpoint_swap}
    \lean{MPSTensor.MPOSymmetry.jointMixedPhysicalSwap,
        MPSTensor.MPOSymmetry.jointMixedEndpointPhysicalSwap}
    \leanok
    \uses{def:mpo_joint_mixed_family}
    Write $\mathcal I$ for the joint alphabet of $(A_0,A_1)$ and
    $\widetilde{\mathcal I}$ for that of the reversed pair $(A_1,A_0)$.
    Define $e:\mathcal I\longrightarrow\widetilde{\mathcal I}$ by
    \begin{align*}
        00(i)     & \longmapsto11(i),     & 11(j) & \longmapsto00(j), \\
        01(x,a,b) & \longmapsto10(x,a,b), &
        10(x,a,b) & \longmapsto01(x,a,b).
    \end{align*}
    Both rectangular indices retain their order. Let
    $s_x:\mathbb C^{D_0^x}\oplus\mathbb C^{D_1^x}
        \longrightarrow\mathbb C^{D_1^x}\oplus\mathbb C^{D_0^x}$
    exchange the virtual summands, and put
    $\Phi_x(X)=s_xXs_x^{-1}$.
    The physical configuration isometry goes from the reversed chain
    space to the original one:
    $(U_Nv)_\sigma=v_{e\circ\sigma}$.
\end{definition}

\begin{theorem}[Reflection of the actual family and joint support]
    \label{thm:mpo_joint_endpoint_swap_support}
    \lean{MPSTensor.MPOSymmetry.jointMixedEndpointBase_swap,
        MPSTensor.MPOSymmetry.bondInterpolationMatrix_reflect_swap,
        MPSTensor.MPOSymmetry.jointMixedEndpointInterpolation_reflect_swap,
        MPSTensor.MPOSymmetry.blockInsertedBoundaryMap_jointMixed_reflect_swap,
        MPSTensor.MPOSymmetry.jointMixedEndpoint_extendedSupport_eq_map_reflect_swap,
        MPSTensor.MPOSymmetry.jointMixedEndpointParentInteraction_eq_conj_reflect_swap}
    \leanok
    \uses{def:mpo_joint_endpoint_swap,def:mpo_joint_mixed_support}
    A tilde denotes the reversed pair. For every $\gamma\in\mathbb R$,
    every block $x$, and every physical letter $\mu$,
    \begin{align}
        \widetilde C_x^{e(\mu)}           & =\Phi_x(C_x^\mu),         &
        \widetilde W_x(1-\gamma)          & =\Phi_x(W_x(\gamma)),     &
        \widetilde A_x^{e(\mu)}(1-\gamma) & =\Phi_x(A_x^\mu(\gamma)).
        \label{eq:mpo_joint_endpoint_swap_letters}
    \end{align}
    Consequently,
    \begin{align}
        U_2\widetilde\Gamma'_{2,1-\gamma}((\Phi_xX_x)_x)
                   & =\Gamma'_{2,\gamma}(X),                \\
        S_\gamma   & =U_2\widetilde S_{1-\gamma},         &
        h'(\gamma) & =U_2\widetilde h'(1-\gamma)U_2^{-1}.
        \label{eq:mpo_joint_endpoint_swap_support}
    \end{align}
    These identities require no spanning or dimension assumptions.
\end{theorem}
\begin{proof}
    \leanok
    \uses{thm:mpo_joint_phase_boundary}
    Inspect the four physical summands. An endpoint letter stays in its
    corresponding virtual corner after both labels are exchanged; a
    rectangular matrix unit keeps its row and column indices. Exchanging
    the two diagonal weights replaces $\gamma$ by $1-\gamma$.
    In each summand of the joint boundary trace, the whole product is
    conjugated by $s_x$, so its trace is unchanged. Every reversed virtual
    boundary is $\Phi_xX_x$ for a unique $X_x$, which proves the support
    equality. Orthogonal projection onto an isometric image is conjugated
    by the same isometry, giving the parent interaction identity.
\end{proof}

\begin{theorem}[Transport of the periodic Hamiltonian and component vectors]
    \label{thm:mpo_joint_endpoint_swap_periodic}
    \lean{MPSTensor.periodicInteractionHamiltonianES_physicalReindex_conj,
        MPSTensor.MPOSymmetry.jointMixedEndpoint_periodic_eq_conj_reflect_swap,
        MPSTensor.MPOSymmetry.mpv_jointMixedEndpointInterpolation_reflect_swap,
        MPSTensor.MPOSymmetry.blockPeriodicMpvMapES_jointMixed_reflect_swap}
    \leanok
    \uses{thm:mpo_joint_endpoint_swap_support}
    For every $N\ge0$ and every real $\gamma$,
    \begin{align}
        H'_N(\gamma) & =U_N\widetilde H'_N(1-\gamma)U_N^{-1}, \\
        U_N\lvert V^{(N)}(\widetilde A_x(1-\gamma))\rangle
                     & =\lvert V^{(N)}(A_x(\gamma))\rangle.
        \label{eq:mpo_joint_endpoint_swap_periodic}
    \end{align}
    The same identity holds for an arbitrary common linear combination
    of the component vectors, with its coefficient vector unchanged.
\end{theorem}
\begin{proof}
    \leanok
    \uses{thm:mpo_joint_endpoint_swap_support}
    Relabeling physical letters commutes with restriction to each cyclic
    window and with extension by the spectator identity. It therefore
    conjugates every translated interaction, and summing gives the first
    identity. This argument includes both oriented terms when $N=2$.
    For a component MPS, induction on the word length conjugates its
    complete matrix product by $s_x$. Trace invariance gives the second
    identity, including the empty word. Summing over $x$ retains every
    cross-label overlap.
\end{proof}

\begin{theorem}[The actual second endpoint and a common bound for both endpoints]
    \label{thm:mpo_joint_both_periodic_endpoints}
    \lean{MPSTensor.MPOSymmetry.jointMixedEndpoint_periodic_groundSpace_one_eq_actual,
        MPSTensor.MPOSymmetry.jointMixedEndpoint_periodic_groundSpace_endpoints_eq_actual,
        MPSTensor.MPOSymmetry.exists_uniform_jointMixedEndpoint_periodic_one_gap,
        MPSTensor.MPOSymmetry.exists_uniform_jointMixedEndpoint_periodic_endpoints_gap}
    \leanok
    \uses{thm:mpo_joint_endpoint_swap_periodic,thm:mpo_joint_periodic_gap,
        thm:martingale_isometric_conjugation_gap}
    Suppose that both endpoint families have simultaneous one-site span
    and $D_1^x>0$ for every block. For every $N\ge2$,
    \begin{align}
        \ker H'_N(1) & =\operatorname{span}_{\mathbb C}
        \{\lvert V^{(N)}(A_x(1))\rangle:x=0,\ldots,r-1\}.
        \label{eq:mpo_joint_right_periodic_kernel}
    \end{align}
    There is a $\delta_1>0$, independent of $N$, such that
    $\|H'_N(1)v\|\ge\delta_1\|v\|$ whenever $v\perp\ker H'_N(1)$.
    Positivity of the first endpoint dimensions is unnecessary here.

    If also $D_0^x>0$ for every block, the actual component-span identity
    holds at both endpoints, and one $\delta_*>0$ satisfies
    \begin{align}
        \|H'_N(\gamma)v\|\ge\delta_*\|v\|
        \quad\text{for }\gamma\in\{0,1\},\ N\ge2,
        \ v\perp\ker H'_N(\gamma).
        \label{eq:mpo_joint_both_periodic_gap}
    \end{align}
    Empty block families and zero physical alphabets are permitted.
    The statement concerns the two endpoints only.
\end{theorem}
\begin{proof}
    \leanok
    \uses{thm:mpo_joint_endpoint_swap_periodic,thm:mpo_joint_periodic_gap,
        thm:martingale_isometric_conjugation_gap}
    Apply Theorem~\ref{thm:mpo_joint_periodic_gap} to $(A_1,A_0)$.
    Its required positive dimensions are precisely $D_1^x$.
    Conjugacy sends the reversed kernel to the original kernel; it also
    sends the reversed component span to the actual span in
    (\ref{eq:mpo_joint_right_periodic_kernel}). If $v$ is orthogonal to the
    original kernel, then $U_N^{-1}v$ is orthogonal to the reversed kernel:
    for every reversed kernel vector $z$,
    $\langle z,U_N^{-1}v\rangle=\langle U_Nz,v\rangle=0$.
    Apply the reversed gap bound and use preservation of both norms.
    Finally choose $\delta_*=\min(\delta_0,\delta_1)$ from the two
    common joint-parent endpoint bounds. No interior parameter enters
    this argument.
\end{proof}
